\documentclass[10pt]{article}
\usepackage[T1]{fontenc}
\usepackage[utf8]{inputenc}
\usepackage[french]{babel}
\usepackage{amsmath,amssymb,mathrsfs,mathpazo}
\usepackage{geometry}
\geometry{a4paper,margin=2cm}
\usepackage{graphicx,array,tabularx,tabularray,colortbl,multirow,multicol,booktabs,enumitem,pifont,ulem,fancybox,diagbox,cancel}
\usepackage{pstricks,pstricks-add,pst-plot,pst-tree,pst-node,pst-text,pst-func,pst-math,pst-bspline,pst-3dplot,pst-eucl,pst-solides3d}
\let\pscalculate\undefined
\usepackage{pst-bezier}
\usepackage{tikz,pgfplots,tkz-tab}
\usetikzlibrary{arrows,arrows.meta,patterns,calc,positioning,trees,decorations.markings,shapes.geometric,shapes.misc,shapes.arrows,intersections,decorations.pathmorphing,decorations.pathreplacing,backgrounds,shadows,babel}
\pgfplotsset{compat=1.18}
\usepgfplotslibrary{fillbetween}
\usepackage[np]{numprint}
\usepackage{listings,makecell,pas-tableur,twoopt}
\let\barre\undefined
\usepackage[pstricks]{tablvar}
\definecolor{darkgreen}{RGB}{0,100,0}
\providecommand{\red}{\color{red}}
\providecommand{\blue}{\color{blue}}
\providecommand{\green}{\color{green}}
\providecommand{\black}{\color{black}}
\providecommand{\white}{\color{white}}
\providecommand{\gray}{\color{gray}}
\providecommand{\cyan}{\color{cyan}}
\providecommand{\magenta}{\color{magenta}}
\providecommand{\yellow}{\color{yellow}}
\providecommand{\violet}{\color{violet}}
\providecommand{\orange}{\color{orange}}
\providecommand{\euro}{\texteuro}
\providecommand{\EUR}{\texteuro}
\providecommand{\ieme}{\textsuperscript{e}}
\providecommand{\degres}{\ensuremath{^\circ}}
\providecommand{\degree}{\ensuremath{^\circ}}
\providecommand{\vv}[1]{\overrightarrow{#1}}
\providecommand{\Coord}[1]{#1}
\providecommand{\decosix}{}
\providecommand{\starredbullet}{\ensuremath{\bullet}}
\pagestyle{empty}
\setlength{\parindent}{0pt}
\usepackage{adjustbox}\DeclareUnicodeCharacter{2236}{:}\DeclareUnicodeCharacter{2190}{\ensuremath{\leftarrow}}
\providecommand{\R}{\mathbb{R}}
\providecommand{\N}{\mathbb{N}}
\providecommand{\Z}{\mathbb{Z}}
\providecommand{\Q}{\mathbb{Q}}
\providecommand{\D}{\mathbb{D}}
\providecommand{\C}{\mathbb{C}}
\providecommand{\e}{{\mathrm{e}}}
\providecommand{\vect}[1]{\overrightarrow{#1}}
\providecommand{\vectt}[1]{\overrightarrow{\mathrm{#1}}}
\providecommand{\vvt}[1]{\overrightarrow{\mathrm{#1}}}
\providecommand{\np}[2][]{{#2~#1}}
\providecommand{\vv}[1]{{\overrightarrow{#1}}}
\renewcommand{\d}{{\mathrm{d}}}
\providecommand{\degres}{{^\circ}}
\providecommand{\degre}{{^\circ}}
\providecommand{\blue}{\color{blue}}
\providecommand{\red}{\color{red}}
\providecommand{\enskip}{\quad}
\providecommand{\llbracket}{\lbrack\!\lbrack}
\providecommand{\rrbracket}{\rbrack\!\rbrack}
\providecommand{\euro}{{\text{€}}}
\providecommand{\no}{{\text{n°}}}
\providecommand{\interval}[2]{{[#1;#2]}}
\providecommand{\intervFO}[2]{{[#1;#2[}}
\providecommand{\dotfill}{\ldots}

\providecommand{\ds}{\displaystyle}
\providecommand{\textding}[1]{#1}
\providecommand{\Oij}{\(\mathrm{O};\vec{\imath},\vec{\jmath}\)}
\providecommand{\Oijk}{\(\mathrm{O};\vec{\imath},\vec{\jmath},\vec{k}\)}
\providecommand{\Ouv}{\(\mathrm{O};\vec{u},\vec{v}\)}
\providecommand{\intervFF}[2]{\ensuremath{[#1;#2]}}
\renewcommand{\euro}{\ifmmode\text{\texteuro}\else\texteuro\fi}
\DeclareUnicodeCharacter{00B0}{\ifmmode^{\circ}\else\textdegree\fi}
\DeclareUnicodeCharacter{20AC}{\ifmmode\text{\texteuro}\else\texteuro\fi}
% Énoncés : Éducation nationale. Transcription et corrigés : APMEP, Denis Vergès et contributeurs. Découpage et classement : CoopMaths, Sébastien Lozano, Éric Elter et contributeurs. Adaptation de la présentation : Maths Explorer.
% CC BY-SA : https://coopmaths.fr/www/about/
% https://coopmaths.fr/alea/static/eam/2026/tex/eam_specialite_2026_09_metropole_3_cor.tex
\begin{document}
\noindent\textit{APMEP / CoopMaths. CC BY-SA. Présentation adaptée par Maths Explorer.}\par\medskip
\emph{Pour $t \in [0~;~5]$, $f(t) = \dfrac{4t+1}{\e^t}$ où $t$ est le temps en heures écoulé
depuis le décollage et $f(t)$ l'altitude de la montgolfière en kilomètres, par rapport au
niveau de la mer.}

\begin{enumerate}
	\item Le décollage correspond à $t = 0$ :
	\[f(0) = \dfrac{4 \times 0 + 1}{\e^0} = \dfrac{1}{1} = 1.\]
	Au décollage, la montgolfière est à 1~km d'altitude : c'est l'altitude du village de
	montagne.

	\item $f$ est un quotient $\dfrac{u}{v}$ avec $u(t) = 4t + 1$, $u'(t) = 4$ et $v(t) = \e^t$, $v'(t) = \e^t$
	($v$ ne s'annule pas). Pour tout réel $t$ de $[0~;~5]$ :
	\[f'(t) = \dfrac{4\e^t - (4t+1)\e^t}{\left(\e^t\right)^2}
	= \dfrac{\e^t\left(4 - 4t - 1\right)}{\e^t \times \e^t}
	= \dfrac{3 - 4t}{\e^t}.\]

	\item Pour tout réel $t$, $\e^t > 0$ : le signe de $f'(t)$ est celui de son numérateur $3 - 4t$.

	Or $3 - 4t \geqslant 0 \iff -4t \geqslant -3 \iff t \leqslant \dfrac{3}{4}$.

	\begin{center}
		\begin{tikzpicture}
			\tkzTabInit[lgt=3.4,espcl=2.8]
			{$t$/0.9,$3-4t$/0.8,$\e^t$/0.8,$f'(t)$/0.8,$f$/2.2}{$0$,$\dfrac{3}{4}$,$5$}
			\tkzTabLine{,+,z,-,}
			\tkzTabLine{,+,,+,}
			\tkzTabLine{,+,z,-,}
			\tkzTabVar{-/$1$,+/$\dfrac{4}{\e^{3/4}}$,-/$\dfrac{21}{\e^5}$}
		\end{tikzpicture}
	\end{center}

avec $f(0)=1$ d'après la question 1 et $f(5)=\frac{4\times 5 + 1}{e^5} = \frac{21}{e^5}$.

	$f'(t) > 0$ sur $\left[0~;~\dfrac{3}{4}\right[$, $f'\left(\dfrac{3}{4}\right) = 0$ et
	$f'(t) < 0$ sur $\left]\dfrac{3}{4}~;~5\right]$. La fonction $f$ est donc croissante sur
	$\left[0~;~\dfrac{3}{4}\right]$ puis décroissante sur $\left[\dfrac{3}{4}~;~5\right]$ : elle
	admet un \textbf{maximum en} $\boldsymbol{t = \dfrac{3}{4}}$, égal à
	$f\left(\dfrac{3}{4}\right) = \dfrac{4}{\e^{3/4}}$.

	\item La montgolfière perd de l'altitude dès que $f$ devient décroissante, c'est-à-dire à partir de $t = \dfrac{3}{4}$~h.
	
Or $\dfrac{3}{4} \times 60 = 45$ : la montgolfière monte pendant 45~minutes après le
	décollage.

10~h~50~min $+$ 45~min $=$ 10~h~95~min $=$ 11~h~35~min.

La montgolfière commence à perdre de l'altitude à \textbf{11~h~35}, soit 45~minutes
	après son décollage.
\end{enumerate}

\end{document}
